% Independent research fragment. Root derivation; independent audits passed.
\section{A terminal-forest improvement to the cut constant}
\label{sec:terminal-forest-cut-constant}

The common-functional lemma gives a much better constant when the local spaces come from a graph.
The following argument replaces a vector-space pigeonhole by a forest with few terminals.

\begin{lemma}[A local intersection in a terminal forest]
\label{lem:local-terminal-intersection}
Let $Q$ be a connected loopless multigraph, and let $S$ be a nonempty proper vertex set.
Write $L_v$ for the span of the edge-coordinate functionals at $v$ in $\mathcal Z(Q)^*$.
Put $U=\sum_{v\in S}L_v$ and $B=|\delta_Q(S)|$.
Let $\mathcal V$ be a set of vertices outside $S$ with $L_v\cap U\ne0$.
Except for at most $2B$ members, $\mathcal V$ can be partitioned into at most $2B$ classes.
Each class has a common nonzero functional in all its spaces $L_v$.
\end{lemma}

\begin{proof}
At most $B$ vertices outside $S$ have a neighbor in $S$.
Exclude those vertices.
Choose a spanning tree in each component of $Q-S$.
For each edge of $\delta_Q(S)$, attach a distinct terminal leaf to its endpoint outside $S$.
Let $T$ be the union of paths between terminal leaves in these trees.
Components with fewer than two terminals contribute nothing.
The forest $T$ has at most $B$ leaves.
It has at most $B$ vertices of degree at least three.
After suppressing its degree-two vertices, it has at most $2B$ edges.
Exclude the members of $\mathcal V$ at branching vertices of $T$.

We show that every remaining member lies internally on one of these edges.
Fix $v\in\mathcal V$ with no neighbor in $S$, and choose $0\ne a\in L_v\cap U$.
Represent $a$ both by a subset $A$ of the star at $v$ and by a subset of edges incident with $S$.
The symmetric difference of these two representing sets annihilates $\mathcal Z(Q)$.
It is therefore a cut of $Q$.
Let $H$ be the component of $Q-S$ containing $v$.
Every component of $H-v$ lies entirely on one side of this cut.
Consequently, $A$ consists of the edges from $v$ to a union of components of $H-v$.
As a functional, $a$ equals the sum of the boundary-terminal coordinates attached to that union.

The sum of all terminal coordinates of $H$ is zero, because $\delta_Q(H)$ is a cut.
Since $a\ne0$, at least two components of $H-v$ contain terminals.
Every spanning tree of $H$ must connect these terminal groups through $v$.
Thus $v$ has degree at least two in $T$.
It has degree exactly two after the branching exceptions have been removed.
There are exactly two terminal-bearing components of $H-v$.
Their terminal sets are the two sides of the corresponding edge of the suppressed forest.
The sum of the terminal coordinates on either side is $a$.
In particular, this common functional is nonzero.

All remaining members on the same suppressed edge give the same terminal bipartition.
They therefore contain the same nonzero functional in their local spaces.
The adjacent and branching exceptions number at most $2B$.
The number of classes is at most $2B$, as required.
\end{proof}

\begin{theorem}[An explicit exponential cut estimate]
\label{thm:terminal-forest-cut-estimate}
Let $Q$ be a connected loopless multigraph with $k$ distinct marked vertices of degree at most $d\ge2$.
Let $Z$ count the marked vertices closed in a uniform even subgraph.
Then
\[
 \mathbb E2^{-Z}\le\frac12+2\exp\!\left(-\frac{k}{2d4^d}\right).
\]
The same bound applies to the probability that a uniform even subgraph is a forest in each of $k$ disjoint cyclic regions.
Each region must have an edge boundary of size at most $d$.
\end{theorem}

\begin{proof}
Use the local spaces $L_v$ from Theorem~\ref{thm:exponential-closed-stars}.
A zero local space gives an expectation at most one half.
If $2^d$ local spaces form a direct sum, that theorem's independence argument also gives an expectation below one half.

Otherwise take a maximal direct-sum family, with vertex set $S$ and size $t<2^d$.
Every remaining local space meets $U=\sum_{v\in S}L_v$ nontrivially.
Suppose first that $k>(1+2d)2^d$.
Then $S$ is nonempty and proper, and $B=|\delta_Q(S)|\le dt$.
Lemma~\ref{lem:local-terminal-intersection} gives a common nonzero functional for at least
\[
 \ell\ge\frac{k-t-2B}{2B}
 \ge\frac{k}{2d2^d}-2
\]
marked vertices.
The common-functional lemma, with $q=1-2^{1-d}\ge1/2$, gives
\[
 \mathbb E2^{-Z}-\frac12
 \le\frac12q^{\lfloor\ell/2\rfloor}
 \le e^{-2^{-d}\ell}
 \le e^{2^{1-d}}\exp\!\left(-\frac{k}{2d4^d}\right)
 \le2\exp\!\left(-\frac{k}{2d4^d}\right).
\]
If $k\le(1+2d)2^d$, the asserted right side exceeds one, so the result is trivial.
The quotient argument in Corollary~\ref{cor:exponential-small-cuts} transfers this estimate to cyclic regions.
\end{proof}

\paragraph{Status and scope.}
The proof has passed independent mathematical audits, including an audit by Fable 5.2.
It is not formalized in Lean.
Its rate is $1/(2d4^d)$, instead of the earlier rate $2^{-O(d2^d)}$.
For an error margin $\varepsilon>0$, more than one half plus that margin forces
\[
 k<2d4^d\ln(2/\varepsilon).
\]
This improves the dependence on $\varepsilon$ from inverse-linear to logarithmic.
